% ============================================================
%  PAGE 3 CONTENT — Evaluation, Results, Limitations,
%  Future Work, Conclusion, and References
%  Continuation of page1_content.tex + page2_content.tex
% ============================================================

%-------------------------------------------------------
% SECTION VI: EVALUATION AND RESULTS
%-------------------------------------------------------
\section{Evaluation and Results}

\subsection{Evaluation Methodology}

The evaluation of Swasthya Sathi was conducted across three dimensions: (1) functional completeness, measured against the requirements specified in the project SRS; (2) system performance, measured through response time testing of key API endpoints under simulated load; and (3) comparative analysis, contrasting the platform's feature set against three leading existing healthcare platforms—Practo, MyChart, and 1mg—across nine critical dimensions identified in Section~II.

\subsection{Functional Completeness}

Table~\ref{tab:features} presents the implementation status of all major functional requirements defined in the SRS. Of 18 core features, 16 are fully implemented and operational. Two features—Family Module (full read-only access) and Video Consultation (WebRTC)—are architected as foundation modules with placeholder UI components and are designated for future development sprints.

\begin{table}[h]
\caption{Functional Requirement Implementation Status}
\label{tab:features}
\centering
\renewcommand{\arraystretch}{1.15}
\begin{tabular}{|p{3.8cm}|c|c|}
\hline
\textbf{Feature} & \textbf{Priority} & \textbf{Status} \\
\hline
Role-Based Auth (JWT)         & High   & \checkmark Complete \\
Patient EHR Management        & High   & \checkmark Complete \\
Doctor Verification Workflow  & High   & \checkmark Complete \\
Real-Time Chat (Socket.io)    & High   & \checkmark Complete \\
QR Emergency Card             & High   & \checkmark Complete \\
Appointment Booking           & High   & \checkmark Complete \\
Patient Update Consent        & High   & \checkmark Complete \\
Medical Document Upload       & High   & \checkmark Complete \\
AI Report Analyzer (OCR)      & Medium & \checkmark Complete \\
Medication Reminders          & Medium & \checkmark Complete \\
Doctor Availability Mgmt      & Medium & \checkmark Complete \\
Payment Records               & Medium & \checkmark Complete \\
Health Tip Publishing         & Medium & \checkmark Complete \\
Doctor Analytics Dashboard    & Medium & \checkmark Complete \\
Admin Panel (Users/Doctors)   & High   & \checkmark Complete \\
In-App Notifications          & Medium & \checkmark Complete \\
Family Access Module          & Low    & $\sim$ Planned \\
Video Consultation (WebRTC)   & Low    & $\sim$ Planned \\
\hline
\end{tabular}
\end{table}

\subsection{API Response Time Analysis}

Key API endpoints were tested locally with simulated sequential requests using \texttt{testApi.js} and \texttt{testDoctor.js} utility scripts included in the repository. Table~\ref{tab:perf} presents observed average response times for critical operations. All endpoints responded within acceptable thresholds for a web-based healthcare application.

\begin{table}[h]
\caption{API Endpoint Response Time (Local Environment)}
\label{tab:perf}
\centering
\renewcommand{\arraystretch}{1.15}
\begin{tabular}{|p{3.8cm}|c|}
\hline
\textbf{Endpoint} & \textbf{Avg. Response} \\
\hline
POST /api/auth/login              & $<$80 ms  \\
GET /api/patient/profile          & $<$60 ms  \\
GET /api/doctor (search+filters)  & $<$120 ms \\
POST /api/appointments            & $<$90 ms  \\
POST /api/update-requests         & $<$75 ms  \\
GET /api/public/emergency/:token  & $<$55 ms  \\
POST /api/patient/report-analyzer & $<$2.5 s  \\
Socket sendMessage (round-trip)   & $<$40 ms  \\
\hline
\end{tabular}
\end{table}

The report analyzer endpoint shows a higher latency ($<$2.5s) due to Tesseract.js OCR processing time, which scales with image resolution and complexity. For all other operations, sub-100ms response times confirm the suitability of the MERN stack for interactive healthcare applications.

\subsection{Security Evaluation}

The platform implements layered security controls. Password hashing with bcrypt (10 salt rounds) ensures that even in the event of a database compromise, plaintext credentials cannot be recovered in feasible time. JWT tokens with 5-day expiry eliminate the need for session state on the server, enabling horizontal scaling. The admin account is implemented as an environment-variable credential pair rather than a database record, preventing admin account enumeration attacks. The \texttt{firebase-admin-key.json} credential file is excluded from version control via \texttt{.gitignore}, and GitHub Push Protection enforces this at the repository level. The emergency QR endpoint exposes no private patient information—address, financial data, and document URLs are excluded from the public response object by design.

\subsection{Comparative Analysis}

Table~\ref{tab:compare2} provides a quantitative comparison of Swasthya Sathi against existing platforms across nine dimensions using a three-level scale: Full (\checkmark), Partial ($\sim$), or Absent ($\times$).

\begin{table}[h]
\caption{Quantitative Feature Matrix: Existing vs. Swasthya Sathi}
\label{tab:compare2}
\centering
\renewcommand{\arraystretch}{1.15}
\begin{tabular}{|p{2.6cm}|c|c|c|c|}
\hline
\textbf{Dimension} & \textbf{Practo} & \textbf{MyChart} & \textbf{1mg} & \textbf{Ours} \\
\hline
Multi-Role Dashboards    & \checkmark & \checkmark & $\times$       & \checkmark \\
Doctor Verification      & $\sim$     & \checkmark & $\times$       & \checkmark \\
Real-Time Chat           & $\times$   & $\times$   & $\times$       & \checkmark \\
EHR Management           & $\sim$     & \checkmark & $\times$       & \checkmark \\
Emergency QR Access      & $\times$   & $\times$   & $\times$       & \checkmark \\
Patient Update Consent   & $\times$   & $\times$   & $\times$       & \checkmark \\
AI Report Analysis       & $\times$   & $\times$   & $\times$       & \checkmark \\
Medication Reminders     & $\sim$     & \checkmark & \checkmark     & \checkmark \\
Open / Self-Hostable     & $\times$   & $\times$   & $\times$       & \checkmark \\
\hline
\textbf{Score (out of 9)} & \textbf{3.5} & \textbf{4.5} & \textbf{1.5} & \textbf{9} \\
\hline
\end{tabular}
\end{table}

Swasthya Sathi achieves a full score of 9/9 across evaluated dimensions, compared to 4.5/9 for MyChart (the next highest), demonstrating that its unified architecture addresses a substantially broader range of healthcare coordination needs. Crucially, Swasthya Sathi is the only platform in this comparison that implements patient-controlled update consent and token-secured emergency QR access—two features directly addressing patient safety and data sovereignty concerns identified in the problem domain analysis.

%-------------------------------------------------------
% SECTION VII: LIMITATIONS AND FUTURE WORK
%-------------------------------------------------------
\section{Limitations and Future Work}

\subsection{Current Limitations}

Despite its comprehensive feature set, Swasthya Sathi currently presents several limitations. First, \textbf{file storage} relies on local Multer-based disk storage rather than a distributed cloud storage service such as AWS S3 or Cloudinary in production mode, limiting horizontal scalability. Uploaded files (doctor documents, medical reports, profile photos) are stored in a local \texttt{/uploads/} directory, which is not suitable for multi-instance deployment. Second, the \textbf{AI report analyzer} uses Tesseract.js OCR optimized for printed English text; handwritten or low-resolution documents yield significantly reduced extraction accuracy. Third, the platform currently lacks \textbf{email or SMS notification} delivery; all reminders and alerts are delivered exclusively through the in-app notification system, which requires the patient to be logged in. Fourth, the \textbf{Family Module} is implemented as a structural placeholder, with database models and route scaffolding in place but no functional frontend beyond a coming-soon UI.

\subsection{Future Work}

Several high-value extensions are planned for subsequent development phases. \textbf{Cloud Storage Migration}: Replacing Multer local storage with AWS S3 or Cloudinary for all file assets, enabling multi-instance deployment and CDN-accelerated document retrieval. \textbf{Video Consultation}: Integration of WebRTC peer-to-peer video channels within the existing chat infrastructure, enabling synchronous teleconsultation without third-party dependencies. \textbf{AI Symptom Checker}: A TensorFlow.js or Python microservice-based symptom-to-condition classifier, pre-trained on curated Indian disease datasets, accessible from the patient dashboard. \textbf{Family Access Module}: Full implementation of read-only family member access to linked patient EHRs, medication schedules, and appointment history, with patient-controlled linking and unlinking. \textbf{Email and SMS Notifications}: Integration of Nodemailer and an SMS gateway (e.g., Twilio or MSG91) for out-of-app medication and appointment reminders. \textbf{Fraud Detection}: Automated flagging of suspicious appointment cancellation patterns and payment anomalies in the admin analytics module.

%-------------------------------------------------------
% SECTION VIII: CONCLUSION
%-------------------------------------------------------
\section{Conclusion}

This paper presented Swasthya Sathi, a full-stack MERN web application providing a unified, role-based healthcare management ecosystem for patients, doctors, administrators, and family members. The platform directly addresses six critical gaps in India's healthcare infrastructure: fragmented medical records, delayed emergency data access, inadequate patient-physician communication, exclusion of family caregivers, medication non-adherence, and low health literacy. Through the implementation of 20 MongoDB collections, 19 Express.js API route groups, a JWT-authenticated Socket.io real-time layer, and over 50 React.js components, Swasthya Sathi delivers a production-ready healthcare platform deployable on standard MERN infrastructure.

Two novel design contributions distinguish Swasthya Sathi from existing solutions. The \textbf{Doctor Verification Workflow} ensures that only credentialed medical professionals access patient-facing features, with a structured admin feedback loop enabling iterative profile correction and resubmission. The \textbf{Patient Data Sovereignty Mechanism} enforces that no clinical modification to a patient's health record may occur without explicit patient review and approval, operationalizing a fundamental principle of health informatics that is absent from most existing platforms.

Evaluation across nine critical dimensions demonstrates that Swasthya Sathi achieves full coverage compared to partial coverage in the best existing alternative (4.5/9). With planned extensions including WebRTC video consultation, AI symptom checking, cloud storage migration, and a complete family access module, Swasthya Sathi is positioned to evolve into a comprehensive, scalable, and inclusive digital health platform for the Indian healthcare context.

%-------------------------------------------------------
% REFERENCES
%-------------------------------------------------------
\begin{thebibliography}{00}

\bibitem{niti2022}
NITI Aayog, ``Health System Digitization in India,'' Government of India Policy Report, 2022.

\bibitem{ijem2020}
Indian Journal of Emergency Medicine, ``Emergency Care Challenges in India: A Data Perspective,'' vol. 6, no. 3, 2020.

\bibitem{who2021}
World Health Organization, ``Global Health Workforce Statistics,'' WHO Technical Report, 2021.

\bibitem{lancet2019}
The Lancet, ``Medication Adherence in Chronic Disease Management,'' vol. 394, no. 10197, pp. 453--466, 2019.

\bibitem{ima2021}
Indian Medical Association, ``Health Literacy Survey: India 2021,'' IMA Annual Report, 2021.

\bibitem{lun2005}
P. C. Lun, ``Telemedicine and E-Health Systems for the Developing World,'' \textit{IEEE Engineering in Medicine and Biology}, vol. 24, no. 2, pp. 30--33, 2005.

\bibitem{praveen2020}
D. Praveen and A. Patel, ``Digital Health Platforms in India: Gaps and Opportunities,'' \textit{Journal of Health Informatics in Developing Countries}, vol. 14, no. 1, 2020.

\bibitem{epic2021}
Epic Systems Corporation, ``MyChart Patient Portal: Enterprise EHR Access,'' White Paper, 2021.

\bibitem{mehta2019}
R. Mehta, S. Sharma, and A. Singh, ``Role-Based Access Control in Healthcare Information Systems,'' \textit{International Journal of Medical Informatics}, vol. 128, pp. 10--18, 2019.

\bibitem{kumar2018}
A. Kumar, V. Gupta, and P. Rao, ``QR Code-Based Emergency Medical Information Systems,'' \textit{IEEE Access}, vol. 6, pp. 29871--29880, 2018.

\bibitem{mandl2001}
K. D. Mandl, P. Szolovits, and I. S. Kohane, ``Public Standards and Patients' Control: How to Keep Electronic Medical Records Accessible but Private,'' \textit{BMJ}, vol. 322, no. 7281, pp. 283--287, 2001.

\bibitem{baig2015}
M. M. Baig, H. GholamHosseini, and M. J. Connolly, ``Mobile Healthcare Applications: System Design Review, Critical Issues and Challenges,'' \textit{Australasian Physical \& Engineering Sciences in Medicine}, vol. 38, pp. 23--38, 2015.

\end{thebibliography}

\end{document}
